\section*{Note de lecture}
\addcontentsline{toc}{section}{Note de lecture}

Ce dossier répond aux versions française et anglaise de l’épreuve de
présélection du \ChallengeName. Il décrit le prototype réellement livré, ses
preuves de fonctionnement et ses limites. Les termes «\,traduire\,» et
«\,contenu inédit\,» n’impliquent pas une génération libre : lorsqu’un sens
financier n’est pas suffisamment couvert, KORA conserve le texte source plutôt
que d’inventer une formulation.

\begin{keybox}[Décision d’ingénierie]
KORA est une extension Chrome autonome qui traduit les textes accessibles du
français ou de l’anglais vers le fulfulde de l’Adamaoua
(\texttt{ISO 639-3 : fub}). Elle ne modifie ni le code source du service
financier, ni les valeurs de ses formulaires. Le moteur, les mémoires, le
correcteur et les modèles linguistiques sont embarqués dans l’extension :
aucun contenu de page ne quitte l’appareil.
\end{keybox}

\subsection*{Accès direct aux livrables}

\begin{tabularx}{\linewidth}{@{}L{32mm}Y@{}}
\toprule
\textbf{Dépôt source} &
\href{\RepoURL}{\RepoURL}\\
\textbf{Vidéo de démonstration} &
\href{\DemoVideoURL}{Voir la vidéo MP4 (2 min 54 s)}\\
\textbf{Extension} &
\code{release/kora-extension.zip} et \code{release/kora-unpacked/}\\
\textbf{Modèle exportable} &
\code{release/kora-fub-neural-model.json}\\
\textbf{Environnement} &
\href{\ChallengeURL}{\ChallengeURL}\\
\bottomrule
\end{tabularx}

\vspace{3mm}
\subsection*{Résumé exécutif}

Les interfaces financières restent souvent proposées uniquement en français
ou en anglais. Cette barrière linguistique devient critique lorsque
l’utilisateur doit distinguer un solde, des frais, un bénéficiaire, une
confirmation ou une alerte de sécurité. KORA apporte une couche de traduction
locale au-dessus d’une page existante, sans hook propriétaire et sans contrat
\code{data-i18n}.

\begin{multicols}{2}
\textbf{Ce que fait KORA}
\begin{itemize}
  \item traduit les libellés FR/EN vers l’Adamawa Fulfulde;
  \item suit les mutations d’une application web après chargement;
  \item couvre un parcours simulé : solde, frais, bénéficiaire,
        récapitulatif, confirmation et reçu;
  \item restaure immédiatement la page source;
  \item conserve les raccourcis, le focus et la structure interactive.
\end{itemize}

\columnbreak
\textbf{Ce que KORA refuse de faire}
\begin{itemize}
  \item modifier un montant, une devise, un nom ou une référence;
  \item traduire la valeur saisie dans un champ ou un \code{select};
  \item appeler une API distante de traduction;
  \item inventer une phrase financière insuffisamment comprise;
  \item prétendre remplacer une certification bancaire ou une validation
        native à grande échelle.
\end{itemize}
\end{multicols}

\subsection*{Fiche d’identité vérifiable}

\begin{tabularx}{\linewidth}{@{}L{46mm}Y@{}}
\toprule
\textbf{Forme du prototype} & Extension Chrome Manifest V3\\
\textbf{Sources prises en charge} & Français et anglais\\
\textbf{Langue nationale cible} & Fulfulde de l’Adamaoua, \code{fub}\\
\textbf{Exécution} & Locale, CPU, sans serveur de traduction\\
\textbf{Permissions} & \code{storage}, \code{activeTab}, \code{scripting}\\
\textbf{Archive courante} & 260\,523 octets, soit environ 254,4 Kio\\
\textbf{Surface terminologique} &
504 concepts financiers dédupliqués; 391 clés FR, 390 clés EN,
113 codes préservés\\
\textbf{Modèle neuronal} &
185\,120 paramètres; contexte 3; embeddings 24; couche tanh 96;
poids int8\\
\textbf{Validation du modèle} &
107 phrases Common Voice tenues hors entraînement;
top-1 14,8\,\%; top-5 35,0\,\%\\
\textbf{Validation linguistique} &
\LinguisticValidator; provenance et statut enregistrés par entrée\\
\bottomrule
\end{tabularx}

\begin{warningbox}[Portée des chiffres]
Les métriques top-1/top-5 mesurent la prédiction du mot suivant en
\texttt{fub}; elles ne sont pas présentées comme un score de traduction.
Le sens FR/EN est déterminé en amont par les mémoires, les gabarits et le
verrou sémantique. Le modèle sert à classer la fluidité de candidats déjà
jugés sûrs.
\end{warningbox}

\clearpage
\section*{Réponse synthétique au cahier des charges}
\addcontentsline{toc}{section}{Réponse synthétique au cahier des charges}

\begin{longtable}{@{}L{40mm}L{34mm}L{78mm}@{}}
\toprule
\textbf{Exigence SAFINUM} & \textbf{Réponse} & \textbf{Preuve dans le prototype}\\
\midrule
\endfirsthead
\toprule
\textbf{Exigence SAFINUM} & \textbf{Réponse} & \textbf{Preuve dans le prototype}\\
\midrule
\endhead
Adaptation linguistique & \pass &
Mémoires bilingues FR/EN vers une sortie commune \code{fub}; popup Fulfulde,
Français et English.\\
Parcours financier & \pass &
Tableau de bord, frais, bénéficiaire, récapitulatif, confirmation, reçu,
erreurs et alertes.\\
Préservation des données & \pass &
Masquage temporaire, restauration caractère pour caractère, contrôle du
nombre d’occurrences et abstention en cas d’échec.\\
Indépendance du service & \pass &
Content script générique sur une page HTTP(S), sans modification du code du
site ni hook dédié.\\
Contenu inédit & \partialpass &
\code{MutationObserver} couvre les nouvelles surfaces DOM; gabarits et
glossaire couvrent les messages dont le sens est connu. Une phrase libre
inconnue reste source par sécurité.\\
Opérations simulées & \pass &
Les tests utilisent des identités, comptes, références et transferts fictifs;
KORA n’initie aucune transaction.\\
Confidentialité & \pass &
Aucun \code{fetch}, WebSocket, télémétrie ou stockage de contenu de page.\\
Accessibilité & \pass &
DOM interactif conservé, focus visible, clavier natif, attribut
\code{lang="fub"}, libellés ARIA traduisibles.\\
Contribution technique propre & \pass &
Moteur, verrou d’intégrité, correction orthographique, modèles linguistiques,
pipeline DOM et suite de tests produits par l’équipe.\\
Documentation et reproductibilité & \pass &
Sources, commandes, architecture, dépendances, limites, licences et matrice de
conformité réunies dans ce dossier et le dépôt.\\
\bottomrule
\end{longtable}

\begin{evidencebox}[Positionnement honnête]
L’exigence «\,adaptation à un contenu inédit\,» possède deux dimensions :
\textit{détecter} tout nouveau texte accessible et \textit{comprendre} son
sens. KORA automatise la première pour les surfaces DOM exposées. Pour la
seconde, il traite les intentions couvertes, les fautes non ambiguës et les
phrases variables; il s’abstient sur un sens absent. Cette politique évite
qu’un gain artificiel de couverture ne crée un défaut majeur financier.
\end{evidencebox}
